\section{Conclusions}\label{Conclusions}

In this paper an effort was made to fit three commonly used machine learning algorithms the \gls{mlp}, \gls{rf} and \gls{svr} models on multivariate currency exchange rate data. This was done in order to assess there applicability for forecasting.

Before the selected data could be used with the chosen models appropriate data preparation techniques had to be applied. As a first variables with missing values had to be removed. Next various combinations of intermediary preprocessing steps including log transformation, differencing and scaling were used to create multiple dataset variations for testing. After that came feature creation consisting of the duplication and application of a offset to each predictor variable. Multiple dataset were created with the offset applied ranging from 1 to 12 months. Following this step was feature selection allowing for the elimination of irrelevant feature variables and decorrelation allowing for the removal of all but one of every set of remaining intercorrelated feature variables. The final step of data preparation was the splitting of the dataset into train and test subsets in a 70:30 proportion. Once these steps were completed the models could be applied to the data. Since each model but the \gls{rf} has specific requirements for optimal model performance not all dataset variations were applied to all models.

All models were fitted on the prepared data using the Optuna hyperparameter optimization framework with the exception of the \gls{rw} model that did not require any fitting. For each model the most suitable data preprocessing variation was selected on the basis of the cumulative average \gls{mse}. Once selected the best values of each result for both metrics between models were compared. On the basis of figure \ref{fig:mse-model-cmp} for the \gls{mse} the best performing models are the \gls{mlp} followed by \gls{svr}, \gls{rf}, \gls{arima} and finally \gls{rw}.

In figure \ref{fig:cumulative-mse} comparing the average \gls{mse} of each best model - dataset combination for each offset a common tendancy is visible and shared by all models. The greater the offset the greater the error. All machine learning models exceeded the performance of the \gls{rw} model with few exceptions. In regards to the \gls{arima} model the results are varied with only the \gls{mlp} exceeding its performance with one exception, while the \gls{rf} and \gls{svr} models exhibitng worse performance in over half of the offsets.

